This study has several limitations. First, the VQC was simulated classically;
therefore, wall-clock time does not represent execution on quantum hardware.
Second, the experiment uses only CartPole-v1 and a small number of seeds
(5 for main 100-episode runs, 3 for the extended 500-episode and DQN variant
runs) due to the submission timeline. Third, the parameter budget is extremely small, which
makes the comparison controlled but also limits the representational capacity
of both agents. In the main experiment, both learned agents perform below a
random policy, so the results characterize a failure regime rather than a
successful policy-learning regime. Fourth, only a shallow VQC (L1) is used in
the main experiment. Deeper circuits may be more expressive, but they also
increase simulation cost and may suffer from barren plateaus~\cite{mcclean2018barren,cerezo2021cost}
in gradient-based optimization. Fifth, the timing comparison between MLP
(GPU) and VQC (CPU) conflates hardware and simulation overheads; the reported
VQC time multipliers should be interpreted as upper bounds, not as precise
quantum-vs-classical scaling.

The results should therefore be interpreted as a small controlled empirical
study, not as evidence for or against quantum advantage in reinforcement
learning.

The fuzzy-teacher distillation experiment has a different status from the DQN
experiments. It demonstrates that small models, including a VQC, can imitate a
useful fuzzy controller when aided by domain-informed fuzzy membership features,
but under raw features the VQC fails even in supervised optimization. Therefore,
the distillation results measure representation quality under structured
features, not DQN learning from reward. The fuzzy membership representation
avoids direct action-recommendation encoding, but it is still a domain-informed
representation and should be reported as such.
